\documentclass[conference]{IEEEtran}
\IEEEoverridecommandlockouts
\usepackage{cite}
\usepackage{amsmath,amssymb,amsfonts}
\usepackage{algorithmic}
\usepackage{graphicx}
\usepackage{textcomp}
\usepackage{xcolor}
\usepackage[utf8]{inputenc}
\usepackage[T1]{fontenc}
\usepackage{hyperref}
\def\BibTeX{{\rm B\kern-.05em{\sc I\kern-.025emB}\kern-.08em
    T\kern-.1667em\lower.7ex\hbox{E}\kern-.125emX}}

\title{MTP-Duplex: Attaching Full-Duplex Turn-Taking to Any LLM with an MTP-Style Protocol Companion}

\author{\IEEEauthorblockN{Penhfel et al.}
\IEEEauthorblockA{\textit{private project — custom-duplex-cascade-pt-br/mtp\_like}}}

\begin{document}
\maketitle

\begin{abstract}
Full-duplex speech dialogue requires the conversational model to perform two
tasks at once: generating answer content \emph{and} controlling turn-taking
(when to speak, when to stop, when to acknowledge). The DuplexCascade approach
trains these two skills into a single model by interleaving answer text with
control tokens. We argue this coupling is unnecessary: a small, dedicated
\emph{protocol companion} — fine-tuned on the same duplex micro-turn data —
can act as a MultiToken-Prediction-style (MTP) head for the control stream,
running alongside an otherwise untouched big model. The companion predicts,
at every conversation event, only the next duplex control tag (or nothing),
while the big model only ever writes answer text. We train a 0.8B companion
(Qwen3.5-0.8B) on the DuplexCascade pt-BR duplex dataset (10k dialogues)
augmented with \emph{barge-in} and \emph{silence-semantics} supervision, and
demonstrate: (i) 97.5\% tag-stream accuracy with perfect F1 on the critical
controller decisions (barge-in interruption and backchannel recognition),
(ii) correct controller behavior on hand-built conversation probes, and
(iii) a working end-to-end browser demo where a stock Qwen3-4B model engages
in barge-in-capable full-duplex conversation, driven entirely by the
companion. The approach is model-agnostic: turn-taking can be attached to any
LLM without fine-tuning it.
\end{abstract}

\begin{IEEEkeywords}
full-duplex dialogue, turn-taking, spoken dialogue systems, MTP,
VAD-free, companion model
\end{IEEEkeywords}

\section{Introduction}
\IEEEPARstart{S}{poken} dialogue systems are moving from turn-based
(``push-to-talk'') interaction to \emph{full-duplex} conversation, where both
parties may speak at any time. The DuplexCascade paper~\cite{duplexcascade}
established that a language model can learn the full-duplex protocol directly:
its training data interleaves user and assistant \emph{micro-turns} and
supervises a set of control tokens (``user is speaking'', ``user finish
speaking'', ``user interruption'', ``user backchannel'', ``user is
thinking'', ``system backchannel'') so that the same model that writes the
answer also emits the control stream. The model must therefore be
fine-tuned — deeply — for the turn-taking skill, and the skill is locked into
that specific checkpoint.

This paper proposes a separation of concerns. Instead of teaching the answer
model the protocol, we train a \emph{small companion model} that predicts
\emph{only} the protocol stream, on behalf of the bigger model, while the
bigger model continues to do what it already does: produce answer text.
This is the same division of labor as MultiToken Prediction (MTP) heads
(e.g., DeepSeek~\cite{deepseek}), where a lightweight module predicts tokens
one step ahead alongside the main model — except that our companion predicts
a different vocabulary: the \emph{duplex cascade protocol}, not the next
answer token. On positions where no control token applies it emits nothing
(an empty string that never enters the context). The companion therefore acts
as the \emph{controller} of the system: it decides turn starts and ends,
barge-ins, backchannels and silences, without any voice-activity detection
(VAD).

Concretely we contribute:

\begin{itemize}
\item a training scheme that reuses the exact DuplexCascade pt-BR duplex
dataset (same micro-turns, same phenomena) but with \emph{tag-only}
supervision at item boundaries, plus two controller-focused augmentations:
\emph{barge-in} records (user words arriving over an open assistant
utterance $\rightarrow$ interruption) and \emph{silence-semantics} records
(user utterance $+$ silence $\rightarrow$ ``finish speaking''), including
minimal single-utterance variants;
\item a 0.8B companion (Qwen3.5-0.8B, text-only) fine-tuned with unsloth
QLoRA on a single RTX 4080 in about two hours, reaching 97.5\% tag-stream
accuracy and perfect F1 on the interruption and backchannel decisions;
\item a runtime \emph{orchestrator} that consumes ASR word events and
silence, queries the companion, and drives any stock LLM — no VAD, no
special tokens in the big model's vocabulary;
\item an end-to-end browser demo in which a stock Qwen3-4B model holds a
barge-in-capable full-duplex conversation in Brazilian Portuguese, entirely
controlled by the companion.
\end{itemize}

\section{Background and Related Work}

\subsection{DuplexCascade}
The DuplexCascade paper~\cite{duplexcascade} constructs a full-duplex speech
dialogue pipeline by (1) converting plain dialogues into interleaved
micro-turns: user turns are split into 1--7-token chunks, each answered by a
system tag such as ``user is speaking''; the final chunk is answered by
``user finish speaking'' plus the first answer chunk; system turns are
delivered in ~10-token bursts separated by ``no voice'' user turns; and six
interaction phenomena (pauses, interruptions, backchannels, thinking pauses,
etc.) are simulated probabilistically; and (2) fine-tuning an LLM on the
resulting stream with masked, weighted cross-entropy so that it learns to
interleave answer text and control tokens. The model's own outputs therefore
\emph{are} the control stream, and the pipeline (ASR $\rightarrow$ LLM
$\rightarrow$ TTS) reacts to them.

The key limitation we address: answer quality and turn-taking skill are
co-trained into one set of weights. Attaching the protocol to a new or
better LLM requires a full re-fine-tune of that model with protocol tokens
in its vocabulary, and the model cannot be used with arbitrary serving
stacks.

\subsection{MultiToken Prediction}
DeepSeek-V3~\cite{deepseek} introduces MTP modules: each MTP layer takes the
previous layer's hidden state and the \emph{next token's embedding} and
predicts the token after next, in parallel with the main model's next-token
prediction. The main model's task is unchanged; the MTP head just looks
ahead. Our companion is analogous in \emph{role} but different in subject:
it looks ahead in the \emph{protocol stream} — it predicts which duplex
control token fires at the next boundary — while the main model continues to
predict answer tokens. The big model is completely untouched.

\subsection{Turn-taking without VAD}
VAD-free turn-taking has been explored via ``word events and silence
timers'' (e.g., in streaming ASR pipelines). DuplexCascade itself is
VAD-free at the LLM level: ``no voice'' markers replace VAD, and the model
learns their meaning. Our companion learns the same semantics from the same
data, but as a standalone decision module: it consumes the ASR word stream
and the assistant's generated text and decides the protocol state. This
keeps the big model's serving stack completely standard.

\section{Method}

\subsection{The protocol vocabulary and the companion's role}
Let $\mathcal{T} = \{\texttt{no\_tag}, \texttt{user\_is\_speaking},
\texttt{user\_finish\_speaking}, \texttt{user\_interruption},
\texttt{user\_backchannel}, \texttt{user\_is\_thinking},
\texttt{system\_backchannel}, \texttt{system\_take\_floor},
\texttt{system\_handover}\}$ be the companion's output vocabulary, where
\texttt{no\_tag} is a label-only symbol (\texttt{$<$no tag$>$}) that never
enters the context. At every conversation event the companion maps the
duplex context (the interleaved user/assistant item stream, rendered in
ChatML) to a distribution over $\mathcal{T}$; a tag fires only if its
probability exceeds a per-tag threshold. Fired tags are appended to the
context as assistant items, exactly like the training stream; \texttt{no\_tag}
means ``emit nothing''.

The last two tags implement the \emph{promptable controller}: when a system
instruction (a \emph{rule}) is present in the context, the companion can fire
\texttt{system\_take\_floor} (the system interrupts the user now — an
instant, sub-second response) or \texttt{system\_handover} (control goes to
the big model, which decides what to do, possibly nothing).

\subsection{Training targets from the DuplexCascade data}
We reuse the pt-BR duplex dataset produced by the companion project
(10k dialogues; the same 50k-source micro-turn pipeline as
\cite{duplexcascade}, with variable-length system chunks). For every item
$i$ of a dialogue we define the target as the tag carried by the
\emph{next} assistant item:

\begin{equation}
\text{target}(i) = \begin{cases}
  \text{special}(\text{item}_{i+1}) & \text{item}_i \text{ is user}\\
  \texttt{no\_tag} & \text{item}_i \text{ is assistant}
\end{cases}
\end{equation}

where \texttt{special} is the control token of the following assistant
micro-turn (or \texttt{no\_tag} when the micro-turn is plain text). Each
dialogue becomes one training sequence with masked labels: cross-entropy is
applied only at item boundaries (the final token of each item). All other
positions are $-100$.

This single transformation gives the companion exactly the supervision the
paper's model received, but as a \emph{classification} of the protocol
state instead of a generation task.

\subsection{Controller augmentations}
The vanilla targets cover the paper's phenomena, but two behaviors are
critical for a controller and under-represented:

\paragraph{Streaming-aware supervision (SAS).} The runtime queries the
companion while user words are still arriving, so for every multi-token
user chunk we add partial-prefix records labeled \texttt{user\_is\_speaking}
(words still arriving), except partials of an interrupting chunk, which are
labeled \texttt{user\_interruption}.

\paragraph{Barge-in (controller) records.} At random points inside every
system turn we render the current assistant item as an \emph{open} item
(partial content, no closing token — speech in progress) followed by
(a) the user's next real words $\rightarrow$ \texttt{user\_interruption}
(plus a partial-prefix variant: the first words of the barge-in),
(b) \texttt{$<$no voice$>$} $\rightarrow$ \texttt{no\_tag} (silence during
speech means keep talking), and (c) a sampled pt-BR backchannel word
$\rightarrow$ \texttt{user\_backchannel} (an acknowledgment is not a
barge-in). This is what lets the companion distinguish ``the user barged
in'' from ``the user just said aham'' from ``the user paused''.

\paragraph{Silence-semantics (turn-end) records.} The runtime inserts a
\texttt{$<$no voice$>$} marker on ASR silence, so the data must define what
silence means in every state: non-final chunk $+$ silence $\rightarrow$
speaking (paused mid-speech); \emph{final} chunk $+$ silence $\rightarrow$
finish speaking (turn complete — the answer must start; the original data
never shows this, since finish always directly follows the final chunk);
partial final chunk $+$ silence $\rightarrow$ speaking; and minimal
single-utterance records $[\texttt{u chunk}, \texttt{$<$no voice$>$}]$
(weight 5) for short single-batch utterances, without which the model never
fires \texttt{finish} on a two-item context.

\paragraph{Promptable-controller (rule) records.} To make the controller
\emph{promptable}, we render a system instruction (the rule) as the first
context item and label the boundary after a trigger chunk with
\texttt{system\_take\_floor} or \texttt{system\_handover}: e.g. the CPF rule
``No CPF, apenas números são permitidos. Se o cliente disser uma letra em vez
de um número, interrompa...'' with a letter chunk after a digit dictation
$\rightarrow$ take\_floor. The synthetic scenarios use the \emph{real
dictation shape} (digit chunks with commas, variable count and order, the
full alphabet as letters) so the companion learns the \emph{concept} (number
vs.\ letter, rule trigger vs.\ unrelated words) instead of memorized words;
negatives (rule present $+$ non-trigger chunk $\rightarrow$ the normal tag)
prevent spurious fires, and a held-out probe confirms that without a rule the
system tags never fire. About half of the regular records additionally get a
generic system instruction (``Gerencie a conversa normalmente...'') so the
base controller behaviors survive the instruction-conditioned training.

\subsection{Loss}
We fine-tune with masked cross-entropy and per-token label weights
(tags carry the paper's relative weights, retuned for the 9-way decision:
finish $=3$, interruption $=6$, backchannel $=3$, thinking $=2$, speaking
$=2$, take\_floor $=5$, handover $=3$, \texttt{no\_tag} $=0.5$), plus a
per-record multiplier (\texttt{loss\_weight}) so the rare controller
patterns dominate the gradient (silence negatives $\times 3$, silence-finish
$\times 3$, barge-in positives $\times 1.5$, rule positives $\times 5$).

Because the companion must keep its \emph{language} and
instruction-following ability (the rule mechanism depends on it), the
content positions use an SDFT-style distill term: forward KL from the frozen
base model over the teacher's top-$k$ support ($k=32$),
\[
\mathcal{L} = \underbrace{\frac{1}{N_{\text{sup}}}\sum_{i\,\in\,\text{bounds}}
w_i\,\text{CE}(\hat{p}_i, y_i)}_{\text{tag boundaries}} \;+\;
\lambda\,\underbrace{\frac{1}{N_{\text{sup}}}\sum_{i\,\in\,\text{content}}
D_{\text{KL}}(p^{\text{base}}_i \,\|\, \hat{p}_i)}_{\text{content
(forward-KL)}},
\]
with $\lambda=0.3$. The teacher is the same weights with the adapters
disabled (no second model, no extra VRAM); this preserves the base model's
language capability while the boundaries learn the protocol.

\subsection{Inference}
At runtime, the companion is queried per event (new ASR words, silence
timer, assistant sentence boundary). Its decision is a distribution over
$\mathcal{T}$; a tag fires only if its probability $\geq \theta_\text{tag}$
and, for \texttt{user\_interruption}, only if the last user item is not the
orchestrator-inserted \texttt{$<$no voice$>$} marker (silence must never
hard-stop the assistant). The context rendering must match the training
distribution exactly: user items are always closed
(\texttt{$<$|im\_end|$>$}), even for partial content; only an in-progress
assistant item is rendered open; the rule, when configured, is the first
\texttt{system} item.

\subsection{Orchestrator}
The orchestrator maintains the duplex context and reacts to events:

\begin{itemize}
\item user words arrive $\rightarrow$ query; \texttt{user\_is\_speaking}
$\rightarrow$ keep waiting (chunk closes when the next words arrive);
\texttt{user\_finish\_speaking} $\rightarrow$ close the user turn and start
the big model; \texttt{user\_interruption} $\rightarrow$ hard stop the
answer and TTS;
\item silence (word-gap timer, no VAD) $\rightarrow$ insert
\texttt{$<$no voice$>$} and query; \texttt{finish} $\rightarrow$ start the
answer; \texttt{no\_tag} $\rightarrow$ keep talking; \texttt{thinking}
$\rightarrow$ wait;
\item assistant sentence boundary $\rightarrow$ query; \texttt{no\_tag}
$\rightarrow$ continue;
\item \texttt{user\_backchannel} $\rightarrow$ keep talking;
\texttt{system\_backchannel} $\rightarrow$ play a canned acknowledgment;
\item \texttt{system\_take\_floor} $\rightarrow$ reset the user's turn and
respond \emph{instantly} with a canned phrase over the local TTS (the big
model's 2--4~s generation would never cut the user off); the rest of the
interrupted utterance is not re-fed (no double fire);
\item \texttt{system\_handover} $\rightarrow$ ask the big model to decide
(respond now, or do nothing) with the rule and the user's words.
\end{itemize}

The big model receives plain text only and is otherwise untouched — except
that it now gets the \emph{conversation history} (previous turns and its own
answers, with protocol tags filtered out), which the companion derives from
the duplex context; a stateless big model otherwise repeats itself (e.g.\ a
bank assistant re-asking for a CPF it already requested).

\section{Experiments}

\subsection{Setup}
\paragraph{Data.} The 10k-dialogue pt-BR duplex dataset (sft pipeline, 2026).
The prep builds $\sim$946k records: 10k main streams, 200k SAS, 155k
barge-in, 553k silence-semantics, 28k prompt-rule scenarios, and the
instruction mix. Split 99.5/0.5 with the same seed as the sft project.

\paragraph{Companion.} \texttt{principled-intelligence/Qwen3.5-0.8B-text-only}
(the Qwen3.5-0.8B language model, 0.75B), QLoRA-style LoRA (r=32, $\alpha$=32)
on all linear layers plus embedding/LM-head fine-tuning of the 10 new special
tokens, bf16, unsloth 2026.9.7 (Qwen3.5 fast path with
\emph{fla}/\emph{causal-conv1d} kernels), 3000 steps, effective batch 16
(per-device batch 1 $\times$ 16 accum — batch 2 exceeds the 16~GB VRAM of the
host's WSL2 setup), lr $10^{-4}$, distill $\lambda=0.3$, $k=32$,
$\sim$5h on a single RTX 4080.

\paragraph{Big model.} Stock \texttt{Qwen/Qwen3-4B-Instruct-2507} for the
end-to-end demo (see \S\ref{sec:discussion} for why Qwen3.5-4B is not the
default).

\subsection{Tag-stream accuracy}
Table~\ref{tab:tags} reports per-tag precision/recall/F1 on the held-out
split (3,222 boundary decisions), computed with restricted argmax over
$\mathcal{T}$.

\begin{table}[htbp]
\centering
\caption{Per-tag results of the trained companion (v6) on the held-out
duplex split.}
\label{tab:tags}
\begin{tabular}{lccc}
\hline
\textbf{tag} & \textbf{P} & \textbf{R} & \textbf{F1}\\
\hline
no tag (emit nothing) & 1.00 & 0.95 & 0.973\\
user is speaking & 0.99 & 0.92 & 0.953\\
user finish speaking & 0.79 & 0.98 & 0.877\\
user interruption & 0.95 & 0.99 & 0.969\\
user backchannel & 1.00 & 0.98 & 0.99\\
user is thinking & 0.90 & 1.00 & 0.949\\
system take floor & 1.00 & 1.00 & 1.00\\
system handover & 1.00 & 1.00 & 1.00\\
\hline
accuracy & \multicolumn{3}{c}{0.956}\\
\hline
\end{tabular}
\end{table}

The two rare but critical controller decisions — interruption and
backchannel recognition — stay near-perfect (F1 0.97/0.99), and the new
promptable-controller tags are perfect on the held-out data. The small
accuracy drop versus v5 (0.975) is the price of the instruction-conditioned
training and the distill term.

\subsection{Controller probes}
We additionally probe the runtime query path directly (guards and
thresholds applied), on hand-built conversations: user speaking
(mid-question) $\rightarrow$ speaking; finished question $+$ silence
$\rightarrow$ finish (fires the answer); assistant mid-utterance $+$ user
barge-in $\rightarrow$ interruption; partial barge-in $\rightarrow$
interruption; ``aham'' during assistant speech $\rightarrow$ backchannel
(not an interruption); silence during assistant speech $\rightarrow$ keep
talking; all correct with probability $\geq 0.83$ at the decision point,
with the remaining ones carried by the no-voice guard / the state machine.

Promptable-controller probes (rule rendered as the system item): CPF rule +
letter in the dictation rhythm $\rightarrow$ take\_floor (1.00); the same
rule with digit chunks only $\rightarrow$ no fire; \emph{no rule} $+$ letter
$\rightarrow$ no fire (the system tags never fire without an instruction);
handover rule + trigger $\rightarrow$ handover (1.00); rule + unrelated words
$\rightarrow$ no fire. In the live demo, the rule detection uses STT digits
normalized to words and per-word chunking; the letter fires at 0.60--0.86
once the dictation rhythm is established.

\subsection{End-to-end}
In the browser demo (a Banco Penha banking assistant persona), a stock
Qwen3-4B answers ``tudo bem?'', is asked for a credit card and replies by
asking for the CPF (with conversation history, it never re-asks), and is
interrupted when the user speaks over it. With the prompt rule configured,
dictating ``1, 2, 3, x, 5'' makes the companion fire
\texttt{system\_take\_floor} on the letter: the system instantly cuts the
user off with ``Apenas números são aceitos no CPF.'' while the companion's
decisions drive the entire exchange; the big model never sees a special
token.

\section{Discussion and Limitations}
\label{sec:discussion}

\paragraph{Why not fine-tune the big model.} The whole point is
composability: the companion is a 0.8B sidecar trained in hours on one GPU;
swapping the big model (a better checkpoint, a GGUF, a different family)
requires no retraining of the turn-taking skill.

\paragraph{Turn-end detection.} The learned turn-end signal is
sentence-final punctuation of the final chunk plus silence. The ASR must
therefore punctuate (faster-whisper does by default); unpunctated transcripts
delay the ``finish'' decision until the next word batch or a second silence.

\paragraph{Qwen3.5-4B as the big model.} The user's target big model,
Qwen3.5-4B, verbalizes its chain of thought (``thinking / Thinking
Process:'' text) even with thinking disabled in the chat template and strong
system directives — the same behavior that led the sft project to reject it.
The demo therefore defaults to Qwen3-4B-Instruct-2507; the companion is
model-agnostic and the 4B remains selectable (e.g., served by llama.cpp).

\paragraph{Barge-in during a closed sentence.} Barge-ins over an
\emph{open} assistant item are detected with probability 1.0; barge-ins over
a closed sentence are detected with $\sim$0.8 probability. The orchestrator
stops the answer on any user words regardless, so the system behaves
correctly even when the tag is marginal.

\paragraph{Data quality lesson.} Two bugs in the data pipeline were found by
eval-driven debugging (a shared token cache corrupting input text across
dialogues, and a duplicated chunk in silence records); fixing them moved
held-out accuracy from 0.70 to 0.98. We recommend always verifying prep
output against the source dialogues.

\paragraph{Prompt-rule detection vs.\ hardcoded words.} The rule mechanism
is learned, not coded: the held-out data shows perfect precision/recall on
the system tags, and probes confirm the model generalizes the letter rule to
unseen letters (a--z) and number words. The one remaining weakness is
rhythm-dependent: a letter that follows just a single digit chunk is a weak
signal (0.17--0.60), while the same letter after two or more digit chunks
fires reliably (0.60--0.86). A data variant with more digit-count diversity
is the planned fix.

\paragraph{System training cost on constrained hardware.} The companion
training must run within 16~GB VRAM under WSL2, where the dxg driver mirrors
GPU memory into host RAM: a per-device batch of 2 (the original setting)
deterministically OOMs the 31~GB VM (15.8~GB VRAM $\rightarrow$ $\sim$31~GB
host anon-RSS). Batch 1 $\times$ 16 accumulation keeps VRAM at
$\sim$11~GB at the same effective batch, at the cost of $\sim$5~h wall time.

\section{Conclusion}
We presented MTP-Duplex: attach full-duplex turn-taking to any LLM by
training a small protocol companion that predicts the DuplexCascade control
stream on behalf of the big model. On the DuplexCascade pt-BR dataset the
companion reaches 95.6\% tag accuracy with near-perfect
interruption/backchannel recognition, and its \emph{promptable controller}
extends the protocol with instruction-conditioned system tags
(take-the-floor and handover) learned under a distill objective that
preserves instruction-following. The companion successfully drives a stock
big model in a live full-duplex browser demo — banking assistant flow with
instant rule interrupts, no VAD, no big-model fine-tuning. The approach
decouples answer quality from turn-taking skill and makes the duplex
protocol a reusable, attachable module.

\begin{thebibliography}{00}
\bibitem{duplexcascade}
J.~Yang, Y.~Fujita, and Y.~Sudo, ``DuplexCascade: full-duplex speech-to-speech
dialogue with VAD-free cascaded ASR--LLM--TTS pipeline and micro-turn
optimization,'' 2026, arXiv:2603.09180. Demo: sbintuitions.github.io/DuplexCascadeDemo.
\bibitem{deepseek}
DeepSeek-AI, ``DeepSeek-V3 technical report,'' 2024, arXiv:2412.19437.
\bibitem{qwen35}
Qwen~Team, ``Qwen3.5: towards native multimodal agents,'' Feb.~2026,
qwen.ai/blog?id=qwen3.5.
\bibitem{unsloth}
Unsloth~AI, ``Unsloth: faster fine-tuning,'' unsloth.ai.
\end{thebibliography}

\end{document}